\section*{Chapitre \ref{ch:b3:mantle-convection} --- Convection mantellique et forces motrices des plaques}

\begin{solution}{exo:b3:mantle-convection:1}
$\mathrm{Ra} = 3500\times9.9\times\num{3e-5}\times1000\times(\num{6.7e5})^3/
(10^{-6}\times\num{5e20}) = \num{6.25e5}$, soit environ 950 fois le \omterm{def:b3:mantle-convection:critical}{nombre de Rayleigh
critique} de 657.5~: la couche convecte vigoureusement.
\end{solution}

\begin{solution}{exo:b3:mantle-convection:2}
$\tau_M = 10^{21}/\num{7.7e10} = \qty{1.30e10}{s}$, soit environ \qty{412}{an}. La
marée, qui dure quelques heures, rencontre un manteau élastique~; la fonte d'une calotte glaciaire,
sur des milliers d'années, et l'enfoncement d'un panneau plongeant, sur des millions, rencontrent un
manteau visqueux.
\end{solution}

\begin{solution}{exo:b3:mantle-convection:3}
$\tau = 4\pi\times10^{21}/(3300\times9.8\lambda)$~: \qty{3.89e11}{s} =
\qty{12.3}{ka} pour \qty{1000}{km}, et \qty{4.1}{ka} pour \qty{3000}{km}. La
charge plus large se rétablit plus vite.
\end{solution}

\begin{solution}{exo:b3:mantle-convection:4}
$g\rho_m\alpha(T_m - T_0)\kappa t = 9.8\times3300\times\num{3.1e-5}\times1450
\times\num{0.8e-6}\times50\times\num{3.156e13} = \qty{1.835e12}{N.m^{-1}}$~; le
facteur entre crochets vaut $1 + 2\times3300\times\num{3.1e-5}\times1450/(\pi\times2270) = 1.0416$,
donc $F_{rp} = \qty{1.91e12}{N.m^{-1}}$, 24 fois plus petite que la traction du panneau plongeant.
\end{solution}

\begin{solution}{exo:b3:mantle-convection:5}
(a) $v = 3300\times9.8\times\num{3e-5}\times300\times(\num{e4})^2/10^{21} =
\qty{2.91e-11}{m.s^{-1}}$~; le temps de montée sur \qty{10}{km} vaut $d/v =
\qty{3.44e14}{s}$, \qty{10.9}{Ma}, le temps de diffusion $d^2/\kappa =
\qty{e14}{s}$, \qty{3.2}{Ma}. Leur rapport vaut $(d^2/\kappa)/(d/v) =
vd/\kappa = \rho g\alpha\Delta Td^3/(\kappa\eta) = \mathrm{Ra} = 0.29$. (b) Elle
se refroidit en premier~: les petites masses ne convectent pas. (c) $\mathrm{Ra} \propto d^3$, donc $d =
10\times(657.5/0.291)^{1/3} = \qty{131}{km}$.
\end{solution}

\begin{solution}{exo:b3:mantle-convection:6}
$2\mathrm{Ra}_c = 1315$ et $\mathrm{Ra}/2\mathrm{Ra}_c = 7604$, dont la racine
cubique vaut 19.66. $\delta = 2891/19.66 = \qty{147}{km}$, $\mathrm{Nu} = 19.66/2 =
9.8$, $u = (10^{-6}/\num{2.891e6})\times19.66^2 = \qty{1.34e-10}{m.s^{-1}} =
\qty{0.42}{cm/an}$, et $q = 9.8\times3.1\times1450/\num{2.891e6} =
\qty{15}{mW.m^{-2}}$.
\end{solution}

\begin{solution}{exo:b3:mantle-convection:7}
$u = 0.09/\num{3.156e7} = \qty{2.85e-9}{m.s^{-1}}$ et $B = 2270\times\num{8e8}
\times\num{2.85e-9} = \qty{5.18e3}{kg.s^{-1}} = \qty{5.2}{Mg/s}$, soit environ
trois cinquièmes de celui de Hawaï. Le flux de chaleur est $1170\times5178/\num{3.1e-5} =
\qty{1.95e11}{W}$, soit \qty{0.2}{TW}, ou \qty{0.42}{\percent} de la
perte thermique de la Terre.
\end{solution}

\begin{solution}{exo:b3:mantle-convection:8}
$v = 2\times30\times9.8\times(\num{4e5})^2/(9\times10^{21}) =
\qty{1.05e-8}{m.s^{-1}} = \qty{0.33}{m/an}$~; pour \qty{2000}{km}, il faut
$\num{2e6}/\num{1.05e-8} = \qty{1.9e14}{s}$, soit environ \qty{6}{Ma}.
\end{solution}

\begin{solution}{exo:b3:mantle-convection:9}
$w = 300e^{-10/4.4} = \qty{30.9}{m}$ et $\dot w = w/\tau = 30.9/4400 =
\qty{7.0}{mm/an}$, soit les deux tiers du taux mesuré près d'Umeå~: un
demi-espace homogène de ce $\tau$ est dans le bon ordre de grandeur, mais le rebond réel est
plus rapide, ou sa dépression plus profonde.
\end{solution}

\begin{solution}{exo:b3:mantle-convection:10}
(a) Avec $w = 2\rho_m\alpha\Delta T\sqrt{\kappa t/\pi}/(\rho_m - \rho_w)$,
\[
  \tfrac12(\rho_m - \rho_w)gw^2 = \frac{2g\rho_m^2\alpha^2\Delta T^2\kappa t}
  {\pi(\rho_m - \rho_w)} = g\rho_m\alpha\Delta T\kappa t\times
  \frac{2\rho_m\alpha\Delta T}{\pi(\rho_m - \rho_w)} .
\] (b) Par parties,
$\int_0^\infty x\operatorname{erfc}x\,\mathrm dx = [\tfrac{x^2}2\operatorname{erfc}
x]_0^\infty + \int_0^\infty\tfrac{x^2}2\,\tfrac2{\sqrt\pi}e^{-x^2}\,\mathrm dx =
0 + \tfrac1{\sqrt\pi}\cdot\tfrac{\sqrt\pi}4 = \tfrac14$. (c) À \qty{100}{Ma},
$\rho_mg\alpha\int z\theta = \qty{3.67e12}{N.m^{-1}}$ et, avec $w =
\qty{3705}{m}$, $\tfrac12(\rho_m - \rho_w)gw^2 = \qty{1.53e11}{N.m^{-1}}$~: le
terme dû à l'eau ne représente que \qty{4}{\percent} de la force.
\end{solution}

\begin{solution}{exo:b3:mantle-convection:11}
$\ln f = 3\ln(\pi^2 + k^2) - 2\ln k$, dont la dérivée $6k/(\pi^2 + k^2) - 2/k$
s'annule lorsque $3k^2 = \pi^2 + k^2$, $k^2 = \pi^2/2$. Alors $f = (3\pi^2/2)^3/
(\pi^2/2) = 27\pi^4/4$. La longueur d'onde $2\pi/k = 2\sqrt2\,d = 2.83\,d$ couvre
une paire de cellules contrarotatives, chacune de largeur $\sqrt2\,d = 1.41\,d$.
\end{solution}

\begin{solution}{exo:b3:mantle-convection:12}
(a) $u = \qty{1.90e-9}{m.s^{-1}}$, $\eta u/h = 10^{20}\times\num{1.90e-9}/
\num{1.5e5} = \qty{1.27}{MPa}$, et sur \qty{6000}{km}, \qty{7.6e12}{N.m^{-1}}.
(b) Le double de la poussée à la dorsale et le sixième de la traction du panneau plongeant~: la \omterm{def:b3:mantle-convection:basal-drag}{traînée basale} peut résister
à la seule poussée à la dorsale, mais pas au panneau plongeant. (c) $\eta uL/h = \num{3.8e12} +
\num{4.6e12}$ donne $u = \num{8.4e12}\times\num{1.5e5}/(\num{6e6}\times10^{20})
= \qty{2.1e-9}{m.s^{-1}} = \qty{6.6}{cm/an}$, soit la vitesse des
plaques océaniques rapides.
\end{solution}

\begin{solution}{pb:b3:mantle-convection:1}
\textbf{1.} $p = 1000\times9.8\times300 = \qty{2.94}{MPa}$.

\textbf{2.} $w = \rho_wh/\rho_m = 300\times1000/3300 = \qty{90.9}{m}$. La croûte
s'enfonce en bloc~; ce qui s'écoule de sous la charge, puis revient, est le
manteau, donc la colonne est compensée par du manteau déplacé en profondeur.

\textbf{3.} $\pi\times(\num{1.25e5})^2 = \qty{4.91e10}{m^2}$, et l'eau
pèse $\num{4.91e10}\times300\times1000 = \qty{1.47e16}{kg}$.

\textbf{4.} Le rivage s'est formé de niveau, lorsque le fond était le plus enfoncé au
centre~; le centre s'est depuis davantage soulevé que le bord, si bien que la ligne
remonte vers le centre, de \qty{77.2}{m} sur \qty{125}{km}~:
\qty{0.62}{m/km}.

\textbf{5.} $\lambda/2 = \qty{250}{km}$, donc $\lambda = \qty{500}{km}$ et $k =
2\pi/\num{5e5} = \qty{1.257e-5}{m^{-1}}$.

\textbf{6.} $1/k = \lambda/2\pi = \qty{79.6}{km}$.

\textbf{7.} $w = 90.9 - \text{soulèvement}$~: 79.4, 69.3, 52.9, 40.4, 26.9 et
\qty{13.7}{m}.

\textbf{8.} $w = w_0e^{-t/\tau}$ donne $\ln w = \ln w_0 - t/\tau$, une droite
de pente $-1/\tau$.

\textbf{9.} $\tau = (14 - 1)/\ln(79.4/13.7) = 13/1.7571 = \qty{7.40}{ka}$.

\textbf{10.} $w(6) = 79.4\,e^{-5/7.40} = \qty{40.4}{m}$, soit un soulèvement de
\qty{50.5}{m}, comme dans le tableau.

\textbf{11.} \qty{13.7}{m} restent~; $\dot w = 13.7/7400 = \qty{1.85}{mm/an}$.

\textbf{12.} La totalité des \qty{13.7}{m}, atteinte de plus en plus lentement.

\textbf{13.} Le taux décroît comme $e^{-t/\tau}$~: $t = 7.40\ln(1.85/0.5) =
7.40\times1.309 = \qty{9.7}{ka}$, soit environ 9700 ans de plus.

\textbf{14.} $\tau = 7399\times\num{3.156e7} = \qty{2.335e11}{s}$, et $\eta =
\tau\rho_mg\lambda/(4\pi) = \num{2.335e11}\times3300\times9.8\times\num{5e5}/
12.566 = \qty{3.0e20}{Pa.s}$.

\textbf{15.} Juste en dessous de la limite inférieure de la fourchette fennoscandienne. L'écoulement
sous une charge de longueur d'onde \qty{500}{km} décroît sur \qty{80}{km}~: le
rebond échantillonne le manteau le plus superficiel et le plus faible, sous et en partie dans ce qui
est, dans la Terre réelle, la lithosphère.

\textbf{16.} $\tau = 4\pi\times\num{3.0e20}/(3300\times9.8\times\num{3e6}) =
\qty{3.9e10}{s} = \qty{1.2}{ka}$. Sous une charge aussi large, l'écoulement atteint
jusqu'à \qty{477}{km} et davantage, vers un manteau inférieur dix fois plus visqueux,
ce qui le ralentit.

\textbf{17.} Elle porte une partie d'une charge étroite par flexion~: la dépression à
l'équilibre est moins profonde et plus large que la valeur isostatique de \qty{90.9}{m}, et une partie du
rebond est élastique et immédiat. Un demi-espace visqueux homogène ajusté au
soulèvement donne alors une viscosité qui mêle les deux effets.

\textbf{18.} $\tau_M = \num{3.0e20}/\num{7.7e10} = \qty{3.9e9}{s}$, soit environ
\qty{124}{an}, soixante fois plus court que $\tau$~: sur des milliers d'années le
manteau se comporte comme un fluide visqueux.

\textbf{19.} $1 - e^{-50/7.40} = 0.9988$~: \qty{99.9}{\percent} de la
dépression d'équilibre a été atteinte, et l'hypothèse est vérifiée.

\textbf{20.} $\eta \propto \lambda$~: diviser $\lambda$ par deux divise $\eta$ par deux, et
$\log_{10}\eta$ diminue de $\log_{10}2 = 0.30$.

\textbf{21.} $\mathrm{Ra} = 3700\times9.9\times\num{3.1e-5}\times1450\times
(\num{6.7e5})^3/(10^{-6}\times\num{3.0e20}) = \num{1.65e6}$.

\textbf{22.} $\num{1.65e6}/657.5 \approx 2500$ fois la valeur critique.

\textbf{23.} $(\mathrm{Ra}/2\mathrm{Ra}_c)^{1/3} = (1254)^{1/3} = 10.78$~: $\delta
= 670/10.78 = \qty{62}{km}$ et $\mathrm{Nu} = 5.4$.

\textbf{24.} $u = (10^{-6}/\num{6.7e5})\times10.78^2 = \qty{1.73e-10}{m.s^{-1}}
= \qty{0.55}{cm/an}$, deux à quatorze fois moins vite que les plaques.

\textbf{25.} $q = 5.4\times3.1\times1450/\num{6.7e5} = \qty{36}{mW.m^{-2}}$. La
couche néglige la chaleur produite dans le manteau et celle venant du
noyau et du manteau inférieur~; et la loi d'échelle abandonne ses constantes d'ordre un,
alors que les plaques, qui se refroidissent comme des demi-espaces à partir des dorsales, perdent bien plus de chaleur sur
le plancher jeune qu'une couche limite moyenne n'en perdrait.

\textbf{26.} $\eta = \qty{3.0e20}{Pa.s}$~: $\log_{10}\eta =
\textbf{\num[text-series-to-math=true]{20.48}}$.
\end{solution}
